\section{存在的问题与困难}

截至目前，本课题仍然存在以下几个方面的问题与困难：

\begin{itemize}
	\item 横向对比实验尚未完全收口。当前已经具备横向实验运行器、部分真实结果和结果沉淀入口，但正式对照统计、统一图表和结论整理尚未最终完成，因此算法收益的量化证明仍需继续补齐。
	\item 项目工作量的定量统计尚未补充。虽然系统演进、算法实现、测试验证和材料沉淀已经形成较大规模的阶段性成果，但目前论文中仍缺少一组统一核实后的数字指标来直接体现工作量，这会影响论文对整体投入规模的呈现效果。
	\item 运行时控制收益仍需通过更清晰的对照方式表达。当前算法主链已经落地，但若后续实验设计和对照基线组织不够清楚，论文容易再次回到“系统功能较多、方法收益不够突出”的叙事风险。
	\item 部分真实工具链仍受外部条件限制。多工具、多阶段系统在实际运行中仍会受到远程服务稳定性、平台环境和资源配置的影响，从而给实验复现和演示过程带来不确定性。
	\item 工程材料向学术表达的转化仍需继续收束。当前设计文档、实现进展、验证报告和图示资源已经较为充分，但要将其进一步压缩、统一并转写为围绕核心算法展开的学术叙述，仍需要持续打磨章节结构和术语口径。
\end{itemize}

总体来看，当前问题主要集中在三个层面：一是实验比较尚未彻底收口，二是工作量与结果仍需进一步量化，三是论文表达还需要继续围绕核心算法统一主线。这些问题会影响论文整体的说服力和完成度，但并不改变“核心方法已经形成、系统基础已经具备”的总体判断。
